\section{Taxonomy Generation Engine}
\label{chap:taxonomy}

\subsection{Algorithm 1: Reasoning-Driven Taxonomy Generation}

The taxonomy builder implements Algorithm~1 from the Simula research paper. Given a
dataset description $y$ and target depth $D$, it produces a set of taxonomies
$\mathcal{T} = \{T_i\}_{i=0}^{K}$ that map the conceptual space.

\begin{figure}[H]
\centering
\begin{tikzpicture}[
    node distance=0.6cm and 1.5cm,
    stepbox/.style={rectangle, draw, rounded corners, fill=llmpurple!10,
                 minimum width=3cm, minimum height=0.7cm, font=\small,
                 text width=2.8cm, align=center},
    arr/.style={-{Stealth[length=2mm]}, thick},
]
    \node[stepbox] (f1) {1. Factor Identification};
    \node[stepbox, below=of f1] (f2) {2. Breadth-First Expansion};
    \node[stepbox, below=of f2] (f3) {3. Best-of-N Proposal};
    \node[stepbox, below=of f3] (f4) {4. Critic Refinement};
    \node[stepbox, below=of f4] (f5) {5. Level Planning};
    
    \draw[arr] (f1) -- (f2);
    \draw[arr] (f2) -- (f3);
    \draw[arr] (f3) -- (f4);
    \draw[arr] (f4) -- (f5);
    \draw[arr, dashed] (f5.east) -- ++(1,0) |- node[right, pos=0.25, font=\tiny]{next level} (f2.east);
\end{tikzpicture}
\caption{Taxonomy generation stages with iterative refinement.}
\label{fig:taxonomy-stages}
\end{figure}

\subsection{Step 1: Factor Identification}

Given schema metadata, the LLM proposes factors of variation:

\begin{lstlisting}[language=Python,caption={Factor identification prompt}]
SYSTEM: You identify the prime factors of variation for a 
        Text-to-SQL training dataset.
        
USER:   Schema context:
        {schema_context}
        
        Identify 3-7 factors that capture the main dimensions
        of queries users might ask. Examples:
        - "query_complexity" (simple, join, aggregation, subquery)
        - "time_range" (point-in-time, period, YTD, MTD)
        - "entity_type" (customer, order, product, supplier)
        
        Output JSON: {"factors": [{"name": "...", "description": "..."}]}
\end{lstlisting}

\subsection{Step 2: Breadth-First Expansion}

Each factor is expanded into a taxonomy tree. For each node, the LLM proposes children:

\begin{lstlisting}[language=Python,caption={Node expansion (simula\_taxonomy\_builder.py, lines 180--220)}]
async def _expand_node(
    self, node: TaxonomyNode, context: str
) -> list[TaxonomyNode]:
    """Expand a single node with Best-of-N proposals."""
    
    # Step 1: Generate N proposals
    proposals = await self.llm.best_of_n(
        prompt=self._build_expansion_prompt(node, context),
        n=self.config.taxonomy.best_of_n,
        temperature=0.8,
    )
    
    # Merge proposals into candidate set
    candidates = self._merge_proposals(proposals)
    
    # Step 2: Critic refinement
    refined = await self._critic_refine(node, candidates, context)
    
    return refined
\end{lstlisting}

\subsection{Step 3: Best-of-N Proposal}

Multiple proposals are generated to increase coverage:

\begin{lstlisting}[language=Python,caption={Best-of-N sampling}]
async def best_of_n(
    self,
    prompt: str,
    n: int = 5,
    temperature: float = 0.8,
) -> list[str]:
    """Generate N diverse completions."""
    tasks = [
        self.generate(prompt, temperature=temperature)
        for _ in range(n)
    ]
    responses = await asyncio.gather(*tasks)
    return [r.content for r in responses]
\end{lstlisting}

\textbf{Known limitation:} All $N$ proposals use the same temperature, which may reduce 
diversity compared to varying temperature across samples. Empirical testing suggests 
this is acceptable for taxonomy generation where structural diversity (prompt variation) 
dominates over sampling diversity. v2 will evaluate temperature annealing ($T \in [0.6, 1.0]$).

\subsection{Step 4: Critic Refinement}

The critic evaluates and refines the proposed children:

\begin{lstlisting}[language=Python,caption={Critic refinement prompt}]
SYSTEM: You are a taxonomy critic. Evaluate the proposed children
        for COMPLETENESS, SOUNDNESS, and SPECIFICITY.
        
USER:   Parent node: {parent.name} - {parent.description}
        Siblings: {siblings}
        Proposed children: {candidates}
        
        Actions you may take:
        - ADD: missing concepts
        - REMOVE: irrelevant or duplicate nodes
        - MERGE: overlapping concepts
        - EDIT: improve descriptions
        
        Output JSON with refined children list.
\end{lstlisting}

\subsection{Step 5: Level Planning}

After completing a level, a plan is generated for the next level:

\begin{lstlisting}[language=Python,caption={Level planning}]
async def _plan_next_level(
    self, current_level_nodes: list[TaxonomyNode]
) -> str:
    """Generate guidance for next level expansion."""
    prompt = f"""
    Current level nodes: {[n.name for n in current_level_nodes]}
    
    Generate a plan for expanding the next level that ensures:
    1. Consistent granularity across all branches
    2. No overlap between sibling expansions
    3. Complete coverage of the concept space
    """
    response = await self.llm.generate(prompt)
    return response.content
\end{lstlisting}

\subsection{Worked Example: Trial Balance Taxonomy}

\begin{table}[H]
\centering
\caption{Sample Taxonomy for Trial Balance Domain}
\label{tab:tb-taxonomy}
\begin{tabular}{llll}
\toprule
\textbf{Level 0} & \textbf{Level 1} & \textbf{Level 2} & \textbf{Level 3} \\
\midrule
query\_type & aggregation & sum & by\_account \\
            &             &     & by\_department \\
            &             & count & distinct\_accounts \\
            & comparison & variance & MoM \\
            &            &          & YoY \\
            & filter & by\_threshold & exceeds\_100M \\
\bottomrule
\end{tabular}
\end{table}

% =============================================================================
% CHAPTER 5: TRAINING DATA GENERATION ENGINE
% =============================================================================
\newpage